\begin{frame}{\textit{Algospeak}}
\textbf{Algospeak} = ensemble des stratégies linguistiques développées pour contourner la modération algorithmique.\\
\textbf{Principe} : remplacer les termes sensibles par des codes, euphémismes, néologismes, émojis. \\
Théorisé par Adam Aleksic, linguiste américain.
\end{frame}

\begin{frame}{Crypto-langages}
    Expression théorisée par Alexandra Saemmer dans \textit{Le parler fransais des Gilles et John. Enquête sur les crypto-langages militants au sein des plateformes} (2019.)\\
    Ce sont des codes collectifs, un langage parallèle lisible par leur communauté.\\
    Exemple: les gilets jaunes.\\
    \vspace{1em}
    Problème: sert aussi à dissimuler des discours dangereux
\end{frame}

\begin{frame}{Crypto-langages}
    Quelques exemples qu'on peut trouver dans l'article d'A. Saemmer:
    \begin{itemize}
        \item macrotte (Macron)
        \item blinder (blindés)
        \item Cast’amere
        \item Gilles et John
    \end{itemize}
\end{frame}